\medskip

Cet exercice est un questionnaire à choix multiples.
Pour chaque question, une seule des quatre propositions est exacte.
Indiquer sur la copie le numéro de la question et la lettre de la proposition choisie.
Aucune justification n'est demandée.

\smallskip

Pour chaque question, une réponse exacte rapporte un point. Une réponse fausse, une réponse multiple ou l'absence de réponse ne rapporte ni n'enlève de point.

\bigskip

\textbf{Question 1 :}

\medskip

On considère la suite numérique $\left(u_n\right)$ définie pour tout $n$ entier naturel par 
\[u_n = \dfrac{1 + 2^n}{3 + 5^n}.\]

Cette suite :

\smallskip

\begin{tabularx}{\linewidth}{X X}
\textbf{a.~~} diverge vers $+\infty$ &\textbf{c.~~} converge vers 0\\
\textbf{b.~~} converge vers $\dfrac25$ &\textbf{d.~~} converge vers $\dfrac13$.
\end{tabularx}

\medskip

\textbf{Question 2 :}

\medskip

Soit $f$ la fonction définie sur $]0~;~ +\infty[$ par $f(x) = x^2\ln x$.

L'expression de la fonction dérivée de $f$ est :

\smallskip

\begin{tabularx}{\linewidth}{X X}
\textbf{a.~~} $f'(x) = 2x\ln x$.&\textbf{c.~~} $f'(x) = 2$.\\
\textbf{b.~~} $f'(x) = x(2\ln x + 1)$.&\textbf{d.~~} $f'(x) = x$.
\end{tabularx}

\medskip

\textbf{Question 3 :}

\medskip

On considère une fonction $h$ définie et continue sur $\R$ dont le tableau de variations est donné ci-dessous:












\begin{center}
{\renewcommand{\arraystretch}{1.3}
\def\esp{\hspace*{3cm}}
\psset{nodesep=3pt,arrowsize=2pt 3}  
$\begin{array}{|c| *3{c}|}
\hline
 x & -\infty   & \esp & +\infty \\
 \hline
  &   &    & \Rnode{max}{+\infty}   \\
\text{ Variations de $h$ } & &  &  \\
 &     \Rnode{min}{-\infty} & & 
\ncline{->}{min}{max}
\rput*(-2,0.8){\Rnode{zero}{0}}
\rput(-2,1.9){\Rnode{alpha}{1}}
\\
\hline
\end{array}$
}
\end{center}

On note $H$ la primitive de $h$ définie sur $\R$ qui s'annule en $0$.

Elle vérifie la propriété :

\begin{tabularx}{\linewidth}{X X}
\textbf{a.~~}$H$ est positive sur $]-\infty~;~0]$.& \textbf{c.~~} $H$ est négative sur $]-\infty~;~1]$.\\
\textbf{b.~~}$H$ est croissante sur $]-\infty~;~1]$.& \textbf{d.~~}$H$ est croissante sur $\R$.
\end{tabularx}



\textbf{Question 4 :}

\medskip

Soit deux réels $a$ et $b$ avec $a < b$.

On considère une fonction $f$ définie, continue, strictement croissante sur l'intervalle
$[a~;~b]$ et qui s'annule en un réel $\alpha$.

Parmi les propositions suivantes, la fonction en langage Python qui permet de donner
une valeur approchée de $\alpha$ à $0,001$ est:

\smallskip

\begin{center}
\begin{tabularx}{\linewidth}{X X}
\textbf{a.~~}&\textbf{c.~~}\\
\begin{tabular}{l}
def racine$(a, b)$:\\
\quad while abs$(b - a)> = 0.001$:\\
\qquad $m = (a + b)/2$ \\
\qquad if $f(m) < 0$ :\\
\qquad \quad $b = m$\\
\qquad  else:\\
\qquad \quad $a = m$ \\
\quad return $m$
\end{tabular}
&
\begin{tabular}{l}
def racine$(a, b)$:\\
\quad $m = (a + b)/2$\\
\quad while abs$(b - a)<=0.001$:\\
\qquad if $f(m)< 0$ :\\
\qquad \quad $a = m$ \\
\qquad else:\\
\qquad \quad $b = m$ \\
\quad return $m$
\end{tabular}\\
\medskip
\textbf{b.~~}&
\medskip\textbf{d.~~}\\
\begin{tabular}{l}
def racine$(a, b)$: \\
\quad $m = (a + b)/2$\\
\quad while abs$(b - a)>= 0.001$:\\
\qquad if $f(m) < 0$:\\
\qquad \quad$a = m$\\
\qquad else:\\
\qquad \quad $b = m$\\
\quad return $m$\\
\end{tabular}
&
\begin{tabular}{l}
def racine $(a, b)$:\\
\quad while abs $(b - a)>= 0.001$:\\
\qquad $m =(a + b)/2$\\
\qquad if $f(m) < 0$:\\
\qquad \quad $a = m$\\
\qquad else :\\
\qquad\quad$b = m$\\
\quad return $m$
\end{tabular}
\end{tabularx}
\end{center}

\medskip

\textbf{Question 5 :}

\medskip

Une urne contient $10$ boules indiscernables au toucher dont $7$ sont bleues et les autres vertes.

On effectue trois tirages successifs avec remise. La probabilité d'obtenir exactement deux boules vertes est :

\smallskip

\begin{tabularx}{\linewidth}{X X}
\textbf{a.~~}$\left(\dfrac{7}{10}\right)^2 \times \dfrac{3}{10}$&\textbf{c.~~} $\ds\binom{10}{2}\left(\dfrac{7}{10}\right)\left(\dfrac{3}{10}\right)^2$\\
\textbf{b.~~}$\left(\dfrac{3}{10}\right)^2$&\textbf{d.~~}$\ds\binom{3}{2}\left(\dfrac{7}{10}\right)\left(\dfrac{3}{10}\right)^2$
\end{tabularx}

\bigskip
